\begin{register}{H}{tselect}{0x7A0}\label{regdesc:TSELECT}
\regfield{(reserved)}{29}{2}{{0x00000000}}%
\regfield{tselect}{2}{0}{{0x0}}%
\reglabel{Reset}\regnewline%
\begin{regdesc}\begin{reglist}
\label{fielddesc:TSELECT}\item [tselect] 配置触发器单元编号 (0-3)。(R/W)
\end{reglist}\end{regdesc}
\end{register}

\begin{register}{H}{tdata1}{0x7A1}\label{regdesc:TDATA1}
\regfield{type}{4}{28}{{0x2}}%
\regfield{dmode}{1}{27}{{0}}%
\regfield{data}{27}{0}{{0x3e00000}}%
\reglabel{Reset}\regnewline%
\begin{regdesc}\begin{reglist}
\label{fielddesc:TDATA1TYPE}\item [type] 表示触发器类型。仅支持匹配类型 (0x2)，此域保留。(RO)
\label{fielddesc:TDATA1DMODE}\item [dmode] 如果某触发器正在被调试器使用，则此域置为 1。\\
0：在调试模式和机器模式下都能写入 tdata1 和 tdata2\\
1：只有在调试模式下才能写入 tdata1 和 tdata2，其他模式下的写操作将被忽略\\
注意：仅支持调试模式下的写操作\\
(R/W)
\label{fielddesc:TDATA1DATA}\item [data] 配置抽象 tdata1 的内容。由于仅支持匹配类型 (0x2) 触发器，此域将始终被解读为 \hyperref[regdesc:MCONTROL]{mcontrol} 的域。(R/W)
\end{reglist}\end{regdesc}
\end{register}

\begin{register}{H}{tdata2}{0x7A2}\label{regdesc:TDATA2}
\regfield{tdata2}{32}{0}{{0x00000000}}%
\reglabel{Reset}\regnewline%
\begin{regdesc}\begin{reglist}
\label{fielddesc:TDATA2}\item [tdata2] 配置抽象 tdata2 的内容。由于仅支持匹配类型 (0x2) 触发器，此域将始终被解读为 \hyperref[regdesc:MADDRESS]{maddress}。(R/W)
\end{reglist}\end{regdesc}
\end{register}

\begin{register}{H}{tdata3}{0x7A3}\label{regdesc:TDATA3}
\regfield{mvalue}{6}{0}{{0x0}}%
\regfield{mselect}{1}{0}{{0x0}}%
\regfield{Reserved}{25}{0}{{0x0}}%
\reglabel{Reset}\regnewline%
\begin{regdesc}\begin{reglist}
\label{fielddesc:MVALUE}\item [mvalue] 与 mselect 一起使用的数据。(R/W)
\label{fielddesc:MSELECT}\item [mselect] 
0：忽略 mvalue\\
1：仅当 mcontext 的低位等于 mvalue 时，此触发器才会匹配。\\
(R/W)
\label{fielddesc:tdata3_rsvd}\item [Reserved] 读取为 0。(RO)
\end{reglist}\end{regdesc}
\end{register}

\begin{register}{H}{tinfo}{0x7A4}\label{regdesc:TINFO}
\regfield{Reserved}{16}{0}{{0x0}}%
\regfield{info}{16}{0}{{0x0}}%
\reglabel{Reset}\regnewline%
\begin{regdesc}\begin{reglist}
\label{fielddesc:tinfo_rsvd}\item [Reserved] 读取为 0。(RO)
\label{fielddesc:INFO}\item [info] 该域的每一位与 tdata1 中列举的每种可能类型一一对应，位 N 对应于类型 N。如果该位被设置，则当前选择的触发器支持该类型。如果当前选择的触发器不存在，则此域为 1。如果类型不可写，则此寄存器可能未实现，在这种情况下读取它会导致非法指令异常。在这种情况下，调试器可以从 tdata1 读取唯一支持的类型。(RO)
\end{reglist}\end{regdesc}
\end{register}

\begin{register}{H}{tcontrol}{0x7A5}\label{regdesc:TCONTROL}
\regfield{(reserved)}{24}{8}{{0x000000}}%
\regfield{mpte}{1}{7}{{0}}%
\regfield{(reserved)}{6}{1}{{0x00}}%
\regfield{mte}{1}{0}{{0}}%
\reglabel{Reset}\regnewline%
\begin{regdesc}\begin{reglist}
\label{fielddesc:TCONTROLMPTE}\item [mpte] 配置是否使能机器模式下前一个触发器。\\
当 CPU 在机器模式下进入异常，\hyperref[fielddesc:TCONTROLMTE]{mte} 的值会自动写入此域。\\
当 CPU 执行 MRET，此域的值会返回 \hyperref[fielddesc:TCONTROLMTE]{mte}，此域变为 0。\\
(R/W)
\label{fielddesc:TCONTROLMTE}\item [mte] 配置是否使能机器模式下触发器。
当 CPU 在机器模式下进入异常，此域的值会自动写入 \hyperref[fielddesc:TCONTROLMPTE]{mpte}，然后此域变为 0，并且 \hyperref[fielddesc:MCONTROLACTION]{action}=0 的触发器被全局禁用。\\
当 CPU 执行 MRET，\hyperref[fielddesc:TCONTROLMPTE]{mpte} 的值会自动返回此域。\\
(R/W)
\end{reglist}\end{regdesc}
\end{register}

\begin{register}{H}{MCONTEXT}{0x7A8}\label{regdesc:MCONTEXT}
\regfield{Reserved}{26}{0}{{0x0}}%
\regfield{context}{6}{0}{{0x0}}%
\reglabel{Reset}\regnewline%
\begin{regdesc}\begin{reglist}
\label{fielddesc:monctext_rsvd}\item [Reserved] 读取为 0。(RO)
\label{fielddesc:CONTEXT}\item [context] 可在 M 模式和调试模式下可写入。机器模式软件可以将上下文编号写入此寄存器，用于设置仅在该上下文中激活触发器。(R/W)
\end{reglist}\end{regdesc}
\end{register}

\begin{register}{H}{mcontrol}{0x7A1}\label{regdesc:MCONTROL}
\regfield{(reserved)}{4}{28}{{0x2}}%
\regfield{dmode}{1}{27}{{0}}%
\regfield{maskmax}{6}{21}{{0x8}}%
\regfield{hit}{1}{20}{{0}}%
\regfield{(reserved)}{1}{19}{{0}}%
\regfield{(timing)}{1}{18}{{0}}%
\regfield{(sizelo)}{2}{16}{{0}}%
\regfield{action}{4}{12}{{0}}%
\regfield{(reserved)}{1}{11}{{0}}%
\regfield{match}{4}{7}{{0}}%
\regfield{m}{1}{6}{{0}}%
\regfield{(reserved)}{2}{4}{{0}}%
\regfield{u}{1}{3}{{0}}%
\regfield{execute}{1}{2}{{0}}%
\regfield{store}{1}{1}{{0}}%
\regfield{load}{1}{0}{{0}}%
\reglabel{Reset}\regnewline%
\begin{regdesc}\begin{reglist}
\label{fielddesc:MCONTROLDMODE}\item [dmode] 与 \hyperref[regdesc:TDATA1]{tdata1} 中的 \hyperref[fielddesc:TDATA1DMODE]{dmode} 相同。(RW*)
\label{fielddesc:MCONTROLMASKMAX}\item [maskmax]
表示当 \hyperref[fielddesc:MCONTROLMATCH]{match} 为 1 时，支持的最大自然二次幂（NAPOT）范围。始终读取为 0x8，这意味着只能屏蔽最低的 8 位。(RO)
\label{fielddesc:MCONTROLHIT}\item [hit] 如果所选触发器先前已触发，则此值为 1。此位需要手动清除。(R/W)
\label{fielddesc:MCONTROLTIMING}\item [timing] 设置此域以启用加载/存储触发器动作的时序。\\
0：此触发器的动作将在触发它的指令执行之前执行，但在所有前面的指令提交之后执行。\\
1：此触发器的动作将在触发它的指令执行之后执行，且应在下一条指令执行之前执行。\\
(R/W)
\label{fielddesc:MCONTROLSIZELO}\item [sizelo] 此域包含访问地址的最低两位。\\
0x0：触发器将尝试匹配任何大小地址的访问\\
0x1：触发器将尝试匹配 8 位存储器访问\\
0x2：触发器将尝试匹配 16 位存储器访问\\
0x3：触发器将尝试匹配 32 位存储器访问\\
(R/W)
\label{fielddesc:MCONTROLACTION}\item [action] 配置所选触发器在触发时执行动作。有效选项为：\\
0x0：引发断点异常\\
0x1：进入调试模式（仅在 dmode=1 时有效）\\
0x2：在触发器 0 命中时启用追踪\\
0x3：在触发器 1 命中时启用追踪\\
0x4：在触发器 2 命中时启用追踪\\
注意：写入无效值将设置为默认值 0x0。\\
(R/W)
\label{fielddesc:MCONTROLMATCH}\item [match] 配置触发器进行数据/指令地址的匹配操作。有效选项为：\\
0x0：严格字节匹配，即与访问中某个字节对应的地址必须严格匹配 \hyperref[regdesc:MADDRESS]{maddress} 的值\\
0x1：NAPOT 匹配，即访问中至少有一个字节处于  \hyperref[regdesc:MADDRESS]{maddress} 中规定的 NAPOT 区域\\
注意：写入超过最大值的数值会被裁剪为最大值 0x1。\\
(R/W)
\item [见下页……]
\end{reglist}\end{regdesc}
\end{register}
\addtocounter{Regfloat}{-1}
\begin{register}{H}{mcontrol}{0x7A1}
\begin{regdesc}\begin{reglist}
\item [接上页……]
\label{fielddesc:MCONTROLM}\item [m] 置位使选定的触发器在机器模式下操作。(R/W)
\label{fielddesc:MCONTROLU}\item [u] 置位使选定的触发器在用户模式下操作。(R/W)
\label{fielddesc:MCONTROLEXECUTE}\item [execute] 置位使选定的触发器在 CPU 执行具有匹配的虚拟地址的指令之前触发。(R/W)
\label{fielddesc:MCONTROLSTORE}\item [store] 置位使选定的触发器在 CPU 执行具有匹配的数据地址的存储器写操作之前触发。(R/W)
\label{fielddesc:MCONTROLLOAD}\item [load] 置位使选定的触发器在 CPU 执行具有匹配的数据地址的存储器读操作之前触发。(R/W)
\end{reglist}\end{regdesc}
\end{register}

\begin{register}{H}{maddress}{0x7A2}\label{regdesc:MADDRESS}
\regfield{maddress}{32}{0}{{0x00000000}}%
\reglabel{Reset}\regnewline%
\begin{regdesc}\begin{reglist}
\label{fielddesc:MADDRESS}\item [maddress] 配置选定的触发器执行匹配操作时使用的地址。当 \hyperref[regdesc:MCONTROL]{mcontrol} 中的 \hyperref[fielddesc:MCONTROLMATCH]{match}=1 时由 NAPOT 解码。(R/W)
\end{reglist}\end{regdesc}
\end{register}
